% Part: lambda-calculus
% Chapter: lambda-definablity
% Section: arithmetical-functions

\documentclass[../../../include/open-logic-section]{subfiles}

\begin{document}

\olfileid[pt]{lam}{rep}{arf}
\olsection{Funções aritméticas \usetoken{S}{lambda definable}}

\begin{prop}
  \ollabel{prop:succ-ld}
  A função sucessor~$\Succ$ é !!{lambda definable}.
\end{prop}

\begin{proof}
Um termo que !!{lambda define}s a função sucessor é
\begin{align*}
  \fn{Succ} & \ident \lambd[a][\lambd[fx][f ({a} f x)]].
  \intertext{De acordo com nossas convenções, esta é uma abreviação de}
  \fn{Succ} & \ident \lambd[a][\lambd[f][\lambd[x][(f (({a} f) x))]]].
  \intertext{$\fn{Succ}$ é uma função que recebe como argumento um
    número~$a$ e resulta em outra função, $\lambd[fx][f ({a}
      f x)]$. Essa função não é, ela própria, um numeral de Church. No entanto,
    se o argumento $a$ for um numeral de Church, ela se reduz a um. Considere:}
   (\lambd[a][\lambd[fx][f ({a} f x)]])\,\num{n} & \redone
   \lambd[fx][f ({\num{n}} f x)].
   \intertext{O termo interno $\num{n}fx$ é um redex, pois
     $\num{n}$ é $\lambd[fx][f^nx]$. Assim, $\num{n}fx \redone f^nx$ e,
     para o termo inteiro, temos}
   \fn{Succ}\, \num n & \red \lambd[fx][f(f^n(x))],
\end{align*}
isto é, $\num{n+1}$.
\end{proof}

\begin{ex}
  Vejamos o que acontece quando aplicamos $\fn{Succ}$ a~$\num{0}$,
  isto é, $\lambd[fx][x]$. Escreveremos os termos por extenso:
  \begin{align*}
    \fn{Succ}\, \num{0} & \ident (\lambd[a][\lambd[f][\lambd[x][(f (({a} f) 
x))]]])(\lambd[f][\lambd[x][x]])\\
    & \redone \lambd[f][\lambd[x][(f (({(\lambd[f][\lambd[x][x]])} f) x))]] \\
    & \redone \lambd[f][\lambd[x][(f ({(\lambd[x][x])} x))]] \\
    & \redone \lambd[f][\lambd[x][(f x)]] \ident \num{1}\\
  \end{align*}
\end{ex}

\begin{prob}
  O termo 
  \begin{align*}
    \fn{Succ}' & \ident \lambd[{n}][\lambd[fx][{n} f (f x)]]
  \end{align*}
  !!{lambda define}s a função sucessor. Explique por quê.
\end{prob}

\begin{prop}
  \ollabel{prop:add-ld}
  A função adição~$\Add$ é !!{lambda definable}.
\end{prop}

\begin{proof}
A adição é !!{lambda define}d pelos termos  
\begin{align*}
  \fn{Add} & \ident \lambd[{a}{b}][\lambd[fx][{a} f ({b} f x)]]
\intertext{ou, alternativamente,}
  \fn{Add}' & \ident \lambd[{a}{b}][{a}\, \fn{Succ} \, {b}].
\intertext{A primeira representação da adição funciona da seguinte maneira: $\fn{Add}$ primeiro
  recebe dois números, ${a}$ e ${b}$. O resultado é uma função que
  recebe $f$ e $x$ e devolve $af(bfx)$. Se $a$ e $b$ forem os numerais de Church
  $\num{n}$ e $\num{m}$, isso se reduz a $f^{n+m}(x)$,
  que é idêntico a $f^{n}(f^{m}(x))$. Ou, passo a passo:}
  (\lambd[{a}{b}][\lambd[fx][{a} f ({b} f x)]])\num n\,\num m & \redone
  \lambd[fx][\num{n}\, f (\num {m}\, f x)] \\
  & \redone \lambd[fx][\num{n}\, f (f^m x)] \\
  & \redone \lambd[fx][f^n (f^m x)] \ident \num {n+m}.
\intertext{A segunda representação da adição, $\fn{Add'}$, funciona
  de outro modo: aplicada a dois numerais de Church $\num{n}$ e
  $\num{m}$,}
\fn{Add}' \num n \,\num m
& \redone \num{n}\, \fn{Succ}\, \num{m}.
\intertext{Mas $\num{n} f x$ sempre se reduz a $f^n(x)$. Portanto,}
  \num{n}\,  \fn{Succ}\, \num{m} & \red \fn{Succ}^n(\num{m}).
\end{align*}
E, como $\fn{Succ}$ !!{lambda define}s a função sucessor, e aplicar a função
sucessor $n$ vezes a~$m$ resulta em $n+m$, isso, por sua vez,
se reduz a~$\num{n+m}$.
\end{proof}

\begin{prop}
  \ollabel{prop:mult-ld}
  A multiplicação é !!{lambda definable} pelo termo
  \[
  \fn{Mult} \ident \lambd[ab][\lambd[fx][a (b f) x]]
  \]
\end{prop}

\begin{proof}
  Para ver como isso funciona, suponha que apliquemos $\fn{Mult}$ aos numerais de Church
  $\num{n}$ e $\num{m}$: $\fn{Mult} \, \num{n} \, \num{m}$ se reduz a
  $\lambd[fx][\num{n}(\num{m}\, f)x]$. O termo $\num{m} f$ define uma
  função que aplica $f$ a seu argumento $m$ vezes. Consequentemente,
  $\num{n} (\num{m} f) x$ aplica a função ``aplicar $f$ $m$ vezes''
  $n$ vezes a~$x$. Em outras palavras, aplicamos $f$ a $x$, $n\cdot
  m$ vezes. Mas o termo normal resultante é precisamente o numeral de Church
  $\num{nm}$.
\end{proof}

\begin{editorial}
Na verdade, podemos simplificar ainda mais esse termo por redução~$\eta$:
\[
  \fn{Mult} \ident \lambd[ab][\lambd[f][a (b f)]].
\]

Mas então primeiro precisamos explicar a redução~$\eta$.
\end{editorial}

\begin{prob}
A multiplicação pode ser !!{lambda define}d pelo termo
\[
  \fn{Mult}' \ident \lambd[ab][a (\fn{Add}\, b) \num{0}].
\]
Explique por que isso funciona.
\end{prob}

A definição da exponenciação como um termo-$\lambd$ é
surpreendentemente simples:
\[
  \fn{Exp} \ident \lambd[be][e b].
\]
O primeiro argumento $b$ é a base, e o segundo,~$e$, é o expoente.
Intuitivamente, $e f$ é $f^e$ segundo nossa codificação dos números. Se isso
for difícil de entender, ainda podemos definir a exponenciação por
multiplicação iterada:
\[
  \fn{Exp}' \ident \lambd[be][e (\fn{Mult}\, b) \num{1}].
\]

O predecessor e a subtração de numerais de Church não são tão simples quanto
poderíamos pensar: exigem a codificação de pares.

\end{document}
